% ===========================================================================
%  ABAC-Charpente — Guide Socratique du Code
%  Comprendre Python et l'EC5 à travers ce logiciel
% ===========================================================================
\documentclass[12pt, a4paper, openany]{book}

% ── Encodage & langue ───────────────────────────────────────────────────────
\usepackage[utf8]{inputenc}
\usepackage[T1]{fontenc}
\usepackage{lmodern}
\usepackage[french]{babel}

% ── Mise en page ────────────────────────────────────────────────────────────
\usepackage[top=2.5cm, bottom=2.5cm, left=3cm, right=2.5cm]{geometry}
\usepackage{fancyhdr}
\pagestyle{fancy}
\fancyhf{}
\fancyhead[LE,RO]{\thepage}
\fancyhead[LO]{\nouppercase{\rightmark}}
\fancyhead[RE]{\nouppercase{\leftmark}}
\renewcommand{\headrulewidth}{0.4pt}

% ── Mathématiques ───────────────────────────────────────────────────────────
\usepackage{amsmath}
\usepackage{amssymb}

% ── Couleurs ────────────────────────────────────────────────────────────────
\usepackage{xcolor}
\definecolor{bleuMaitre}{RGB}{13, 71, 161}
\definecolor{bleuMaitreLight}{RGB}{227, 242, 253}
\definecolor{ambreApprenti}{RGB}{121, 85, 72}
\definecolor{ambreApprentiLight}{RGB}{251, 233, 231}
\definecolor{vertImportant}{RGB}{27, 94, 32}
\definecolor{vertImportantLight}{RGB}{232, 245, 233}
\definecolor{greyCodeLight}{RGB}{250, 250, 250}
\definecolor{orangeDefinition}{RGB}{230, 81, 0}
\definecolor{orangeDefinitionLight}{RGB}{255, 243, 224}

% ── Boîtes ──────────────────────────────────────────────────────────────────
\usepackage[most]{tcolorbox}

\newtcolorbox{maitre}[1][]{
  enhanced,
  colback=bleuMaitreLight,
  colframe=bleuMaitre,
  fonttitle=\bfseries\color{white},
  title={\textbf{Maître}},
  attach boxed title to top left={yshift=-2mm, xshift=5mm},
  boxed title style={colback=bleuMaitre, colframe=bleuMaitre},
  left=5mm, right=5mm, top=5mm, bottom=3mm,
  #1
}

\newtcolorbox{apprenti}[1][]{
  enhanced,
  colback=ambreApprentiLight,
  colframe=ambreApprenti,
  fonttitle=\bfseries\color{white},
  title={\textbf{Apprenti}},
  attach boxed title to top left={yshift=-2mm, xshift=5mm},
  boxed title style={colback=ambreApprenti, colframe=ambreApprenti},
  left=5mm, right=5mm, top=5mm, bottom=3mm,
  #1
}

\newtcolorbox{important}[1][]{
  enhanced,
  colback=vertImportantLight,
  colframe=vertImportant,
  fonttitle=\bfseries\color{white},
  title={Point cl\'e},
  attach boxed title to top left={yshift=-2mm, xshift=5mm},
  boxed title style={colback=vertImportant},
  left=5mm, right=5mm, top=5mm, bottom=3mm,
  #1
}

\newtcolorbox{definition}[1][]{
  enhanced,
  colback=orangeDefinitionLight,
  colframe=orangeDefinition,
  fonttitle=\bfseries\color{white},
  title={D\'efinition},
  attach boxed title to top left={yshift=-2mm, xshift=5mm},
  boxed title style={colback=orangeDefinition},
  left=5mm, right=5mm, top=5mm, bottom=3mm,
  #1
}

% ── Code ────────────────────────────────────────────────────────────────────
\usepackage{listings}

\lstdefinestyle{python}{
  language=Python,
  backgroundcolor=\color{greyCodeLight},
  basicstyle=\ttfamily\footnotesize,
  keywordstyle=\color{blue!80!black}\bfseries,
  stringstyle=\color{orange!80!black},
  commentstyle=\color{gray!70}\itshape,
  numberstyle=\tiny\color{gray},
  numbers=left,
  breaklines=true,
  showstringspaces=false,
  frame=single,
  rulecolor=\color{gray!40},
  xleftmargin=2em,
  framexleftmargin=2em,
  literate=%
    {à}{{\`a}}1   {â}{{\^a}}1
    {é}{{\'e}}1   {è}{{\`e}}1   {ê}{{\^e}}1   {ë}{{\"e}}1
    {î}{{\^i}}1   {ï}{{\"i}}1
    {ô}{{\^o}}1   {ù}{{\`u}}1   {û}{{\^u}}1   {ü}{{\"u}}1
    {ç}{{\c c}}1
    {É}{{\'E}}1   {È}{{\`E}}1   {Ê}{{\^E}}1
    {À}{{\`A}}1   {Â}{{\^A}}1
    {Î}{{\^I}}1   {Û}{{\^U}}1   {Ç}{{\c C}}1
    {œ}{{\oe}}1   {æ}{{\ae}}1
    {²}{{\textsuperscript{2}}}1
    {³}{{\textsuperscript{3}}}1
    {⁴}{{\textsuperscript{4}}}1
    {§}{{\S}}1
    {×}{{$\times$}}1
    {→}{{$\rightarrow$}}1
    {✓}{{$\checkmark$}}1
    {σ}{{$\sigma$}}1
    {—}{{\textemdash}}1
}

\lstdefinestyle{toml}{
  basicstyle=\ttfamily\footnotesize,
  backgroundcolor=\color{gray!8},
  commentstyle=\color{gray}\itshape,
  frame=single,
  rulecolor=\color{gray!40},
  numbers=left,
  numberstyle=\tiny\color{gray},
  breaklines=true,
  xleftmargin=2em,
  framexleftmargin=2em,
  literate=%
    {à}{{\`a}}1  {â}{{\^a}}1
    {é}{{\'e}}1  {è}{{\`e}}1  {ê}{{\^e}}1
    {î}{{\^i}}1  {ô}{{\^o}}1
    {ù}{{\`u}}1  {û}{{\^u}}1  {ç}{{\c c}}1
    {É}{{\'E}}1  {È}{{\`E}}1
    {À}{{\`A}}1  {Â}{{\^A}}1
    {²}{{\textsuperscript{2}}}1
    {³}{{\textsuperscript{3}}}1
    {⁴}{{\textsuperscript{4}}}1
    {§}{{\S}}1
    {│}{{|}}1
    {├}{{|}}1
    {└}{{|}}1
    {─}{{-}}1
}

% ── Tableaux ────────────────────────────────────────────────────────────────
\usepackage{booktabs}
\usepackage{array}

% ── Graphiques & schémas ────────────────────────────────────────────────────
\usepackage{tikz}
\usetikzlibrary{shapes.geometric, arrows, positioning}

% ── Liens (chargé EN DERNIER pour compatibilité) ─────────────────────────────
\usepackage[colorlinks=true, linkcolor=bleuMaitre, urlcolor=bleuMaitre,
            citecolor=bleuMaitre]{hyperref}

% ── Divers ──────────────────────────────────────────────────────────────────
\usepackage{microtype}
\usepackage{parskip}
\setlength{\parskip}{6pt}

% ── Commandes personnalisées ─────────────────────────────────────────────────
\newcommand{\py}[1]{\texttt{#1}}
\newcommand{\fichier}[1]{\texttt{\textit{#1}}}

% ===========================================================================
\begin{document}
% ===========================================================================

% ── Page de titre ───────────────────────────────────────────────────────────
\begin{titlepage}
\centering
\vspace*{1.5cm}
{\fontsize{28}{34}\bfseries Comprendre ABAC-Charpente}\\[0.6cm]
{\LARGE Guide Socratique du Code}\\[1.5cm]
\rule{\textwidth}{1.5pt}\\[0.5cm]
{\large De Python aux v\'erifications Eurocode 5,\\
expliqu\'e pas \`a pas comme si tu ne savais rien}\\[0.5cm]
\rule{\textwidth}{1.5pt}\\[2cm]
{\large Avril 2026}\\[3cm]

\begin{tcolorbox}[
  colback=bleuMaitreLight, colframe=bleuMaitre,
  width=0.85\textwidth, center
]
\centering\large\itshape
\guillemotleft\, Je ne sais qu'une chose,\\
c'est que je ne sais rien.\,\guillemotright\\[4pt]
{\normalsize --- Socrate}
\end{tcolorbox}

\vfill
{\small Document d'usage personnel --- 2026}
\end{titlepage}

\tableofcontents
\newpage

% ===========================================================================
\chapter*{Avant-propos}
\addcontentsline{toc}{chapter}{Avant-propos}
% ===========================================================================

Ce document est un guide p\'edagogique du logiciel \textbf{ABAC-Charpente}
(\emph{Abaque de Calcul pour Charpentes en Bois}), un outil de calcul
automatique de port\'ees admissibles pour les structures bois,
conform\'ement aux normes europ\'eennes EN~1990 et EN~1995 (Eurocode~5).

La m\'ethode p\'edagogique choisie est le \textbf{dialogue socratique}~: plut\^ot
qu'asséner des d\'efinitions, le Ma\^itre pose des questions qui am\`enent
l'Apprenti \`a d\'ecouvrir les concepts par lui-m\^eme. Tu trouveras tout au
long de ce document des \'echanges entre ces deux personnages~:

\begin{maitre}
Les questions du Ma\^itre t'invitent \`a r\'efl\'echir avant que la r\'eponse
ne soit donn\'ee.
\end{maitre}

\begin{apprenti}
Les r\'eponses de l'Apprenti --- parfois h\'esitantes --- mod\'elisent
le cheminement naturel d'un ing\'enieur qui d\'ecouvre Python.
\end{apprenti}

\begin{important}
Les encadr\'es verts r\'esument les points cl\'es \`a retenir absolument.
\end{important}

\begin{definition}
Les encadr\'es oranges d\'efinissent les termes techniques importants
(Python ou Eurocode).
\end{definition}

\textbf{Comment lire ce document~?}
Tu peux le parcourir lin\'eairement pour une compr\'ehension compl\`ete.
Les exemples de code sont syst\'ematiquement tir\'es de l'application r\'eelle.
Les num\'eros de fichiers mentionn\'es correspondent \`a la structure du projet telle
qu'elle existe sur le disque.

% ===========================================================================
\chapter{Python~: les briques de base}
% ===========================================================================

Avant d'entrer dans le c\oe{}ur du logiciel, il faut parler le m\^eme
langage. Ce chapitre t'apprend Python \`a travers les yeux d'un charpentier.

% ─────────────────────────────────────────────────────────────────────────────
\section{Les variables~: nommer les choses}
% ─────────────────────────────────────────────────────────────────────────────

\begin{maitre}
Quand tu dessines un plan de charpente, tu notes des dimensions.
Comment tu notes la largeur d'une section de bois~?
\end{maitre}

\begin{apprenti}
Je note \guillemotleft\,b = 63 mm\,\guillemotright.
\end{apprenti}

\begin{maitre}
Exactement. En Python, c'est presque identique~:
\end{maitre}

\begin{lstlisting}[style=python, caption={D\'eclaration de variables en Python}]
b_mm = 63.0      # largeur de section, en millimètres
h_mm = 175.0     # hauteur de section, en millimètres
L_max_m = 6.0    # portée maximale, en mètres
\end{lstlisting}

\begin{definition}
Une \textbf{variable} est une bo\^ite dans la m\'emoire de l'ordinateur,
\`a laquelle on donne un nom. Le signe \py{=} ne veut pas dire
\guillemotleft\,est \'egal\,\guillemotright{} au sens math\'ematique,
mais \guillemotleft\,re\c{c}oit la valeur\,\guillemotright.
Autrement dit~: \py{b\_mm = 63.0} se lit
\guillemotleft\,la bo\^ite nomm\'ee \py{b\_mm} re\c{c}oit la valeur
\py{63.0}\,\guillemotright.
\end{definition}

Tu remarques que les noms de variables dans ce code portent l'unit\'e
(\py{\_mm}, \py{\_m}, \py{\_MPa}, \py{\_kNm}).
C'est une convention d\'elib\'er\'ee du projet, appel\'ee \textbf{Principe~VI},
pour \'eviter de confondre m\`etres et millim\`etres ---
une erreur qui a caus\'e plus d'un probl\`eme structurel dans l'histoire
de l'ing\'enierie.

\begin{maitre}
Et si j'\'ecris \py{b\_mm = 63} au lieu de \py{b\_mm = 63.0},
c'est la m\^eme chose~?
\end{maitre}

\begin{apprenti}
Je pense que oui, 63 et 63{,}0 c'est pareil, non~?
\end{apprenti}

\begin{maitre}
Presque. Mais Python distingue les types de donn\'ees~: \py{63} est un
entier (\py{int}), et \py{63.0} est un d\'ecimal (\py{float}).
En pratique, pour les grandeurs physiques, on utilise toujours \py{float}
car les calculs comme $\sigma = M / W$ donnent des r\'esultats d\'ecimaux.
\end{maitre}

% ─────────────────────────────────────────────────────────────────────────────
\section{Les types de donn\'ees}
% ─────────────────────────────────────────────────────────────────────────────

Python distingue plusieurs \emph{types} de donn\'ees~:

\begin{center}
\begin{tabular}{llll}
\toprule
\textbf{Type Python} & \textbf{Ce que c'est}    & \textbf{Exemple dans le projet}\\
\midrule
\py{int}             & Nombre entier             & \py{classe\_service = 1}\\
\py{float}           & Nombre d\'ecimal           & \py{f\_m\_k\_MPa = 24.0}\\
\py{str}             & Texte (cha\^ine)           & \py{classe\_resistance = "C24"}\\
\py{bool}            & Vrai ou Faux              & \py{double\_flexion = True}\\
\py{list}            & Liste d'\'el\'ements       & \py{pentes\_deg = [4, 15, 30, 45]}\\
\py{dict}            & Table cl\'e\,$\rightarrow$\,valeur   & \py{\{"b\_mm": 63, "h\_mm": 175\}}\\
\bottomrule
\end{tabular}
\end{center}

\begin{lstlisting}[style=python, caption={Types de données illustrés par des grandeurs charpente}]
# float : grandeurs physiques avec unité dans le nom
b_mm = 63.0
f_m_k_MPa = 24.0        # résistance caractéristique en flexion (MPa)

# str : identifiants et classifications
classe_resistance = "C24"
type_poutre = "Panne"
usage = "TOITURE_INACC"

# bool : options oui/non
double_flexion = True
second_oeuvre = False

# list : plusieurs valeurs à tester (expansion cartésienne)
pentes_deg = [4, 15, 30, 45]      # 4 pentes à calculer

# dict : propriétés regroupées en une seule structure
section = {
    "b_mm": 63.0,
    "h_mm": 175.0,
    "classe": "C24"
}
\end{lstlisting}

\begin{maitre}
Pourquoi avoir une liste de pentes \py{[4, 15, 30, 45]} plut\^ot qu'une
seule valeur~?
\end{maitre}

\begin{apprenti}
Parce que le logiciel doit calculer pour plusieurs pentes de toiture
diff\'erentes~?
\end{apprenti}

\begin{maitre}
Exactement. Et c'est l\`a qu'intervient l'un des concepts les plus puissants
de ce logiciel~: on peut donner des listes de param\`etres, et le programme
g\'en\`ere automatiquement toutes les combinaisons. Mais n'anticipons pas~---
voyons d'abord les fonctions.
\end{maitre}

% ─────────────────────────────────────────────────────────────────────────────
\section{Les fonctions~: des recettes de calcul r\'eutilisables}
% ─────────────────────────────────────────────────────────────────────────────

\begin{maitre}
Dans ton travail d'ing\'enieur, tu appliques souvent la m\^eme formule
\`a des donn\'ees diff\'erentes. Par exemple, le moment d'inertie d'une
section rectangulaire~:
$I = bh^3/12$.
Tu la recalcules \`a chaque fois depuis z\'ero~?
\end{maitre}

\begin{apprenti}
Non, je la note une fois dans mon tableur et je change juste les valeurs
de $b$ et $h$.
\end{apprenti}

\begin{maitre}
En Python, on fait pareil avec les \textbf{fonctions}.
\end{maitre}

\begin{lstlisting}[style=python, caption={Définition et appel d'une fonction Python}]
# --- Définition : on écrit la "recette" une seule fois ---
def calculer_inertie(b_mm: float, h_mm: float) -> float:
    """Moment d'inertie en cm4 d'une section rectangulaire b x h."""
    I_mm4 = b_mm * h_mm**3 / 12.0    # mm4  (**3 = puissance 3)
    I_cm4 = I_mm4 / 1e4              # conversion mm4 -> cm4
    return I_cm4                      # renvoie le résultat

# --- Appel : on "invoque" la recette avec des valeurs ---
I_A = calculer_inertie(63.0, 175.0)    # -> 28013 cm4
I_B = calculer_inertie(75.0, 225.0)    # -> 71367 cm4
\end{lstlisting}

D\'ecortiquons ce code ligne par ligne~:

\begin{itemize}
  \item \py{def} est le mot-cl\'e Python qui signifie
        \guillemotleft\,je d\'efinis une fonction\,\guillemotright.
  \item \py{calculer\_inertie} est le nom qu'on lui donne
        (en fran\c{c}ais, conforme au Principe~IX du projet).
  \item \py{b\_mm: float, h\_mm: float} sont les \textbf{param\`etres}
        (les ingr\'edients de la recette). Le \py{: float} indique le type
        attendu~: c'est une annotation de type, facultative mais tr\`es lisible.
  \item \py{-> float} indique ce que la fonction retourne.
  \item Les triples guillemets \py{"""..."""} sont la \textbf{docstring}~:
        une documentation int\'egr\'ee, lisible par les outils de d\'eveloppement.
  \item \py{return} renvoie le r\'esultat au code appelant.
\end{itemize}

\begin{important}
L'indentation (les espaces en d\'ebut de ligne) est \textbf{obligatoire}
en Python. Elle d\'elimite le contenu de la fonction. Si tu oublies de
l'indenter, Python l\`eve une erreur imm\'ediatement.
\end{important}

Voici la vraie fonction de calcul de section dans le code~:

\begin{lstlisting}[style=python, caption={Extrait de \fichier{ec5/proprietes.py} --- calcul des propriétés de section}]
def calculer_section(b_mm: float, h_mm: float) -> dict:
    """Calcule les propriétés géométriques d'une section rectangulaire.

    Retourne un dictionnaire avec :
        A_cm2   : aire en cm²
        I_cm4   : inertie axe fort en cm⁴
        W_cm3   : module de résistance axe fort en cm³
        I_z_cm4 : inertie axe faible en cm⁴
        W_z_cm3 : module de résistance axe faible en cm³
    """
    A_mm2   = b_mm * h_mm
    I_mm4   = b_mm * h_mm**3 / 12.0
    W_mm3   = b_mm * h_mm**2 / 6.0
    I_z_mm4 = h_mm * b_mm**3 / 12.0   # axe faible : b et h sont inversés
    W_z_mm3 = h_mm * b_mm**2 / 6.0

    return {
        "A_cm2":   A_mm2  / 100.0,
        "I_cm4":   I_mm4  / 1e4,
        "W_cm3":   W_mm3  / 1e3,
        "I_z_cm4": I_z_mm4 / 1e4,
        "W_z_cm3": W_z_mm3 / 1e3,
    }
\end{lstlisting}

\begin{maitre}
Pourquoi la fonction retourne un \py{dict} plut\^ot que de retourner
\py{A\_cm2} directement~?
\end{maitre}

\begin{apprenti}
Parce qu'on veut r\'ecup\'erer plusieurs valeurs \`a la fois~?
\end{apprenti}

\begin{maitre}
Exactement. En Python, une fonction ne peut retourner qu'une seule
\guillemotleft\,chose\,\guillemotright, mais cette chose peut \^etre un
dictionnaire qui en contient plusieurs. C'est un idiome tr\`es courant
dans ce code~: toutes les fonctions de v\'erification EC5 retournent
un dictionnaire de valeurs interm\'ediaires.
\end{maitre}

% ─────────────────────────────────────────────────────────────────────────────
\section{Les listes et les boucles~: travailler en s\'erie}
% ─────────────────────────────────────────────────────────────────────────────

\begin{maitre}
Imaginons que tu doives v\'erifier une poutre pour des port\'ees de
2\,m, 2,5\,m, 3\,m\ldots{} jusqu'\`a 6\,m. Comment tu ferais manuellement~?
\end{maitre}

\begin{apprenti}
Je ferais le calcul pour 2\,m, puis pour 2,5\,m, puis pour 3\,m\ldots{}
Ca fait beaucoup de lignes.
\end{apprenti}

\begin{maitre}
Python te permet de l'\'ecrire une seule fois et de le r\'ep\'eter
automatiquement avec une \textbf{boucle}.
\end{maitre}

\begin{lstlisting}[style=python, caption={Boucle sur une liste de portées}]
import numpy as np

# Génération d'un tableau de portées : de 2.0 à 6.0, par pas de 0.5
longueurs_m = np.arange(2.0, 6.5, 0.5)
# Résultat : array([2.0, 2.5, 3.0, 3.5, 4.0, 4.5, 5.0, 5.5, 6.0])

# Boucle : on répète le bloc indenté pour chaque valeur de L
for L in longueurs_m:
    w_inst = 5 * 1.2 * L**4 / (384 * 11000 * 2800)  # formule simplifiée
    print(f"L = {L:.1f} m  ->  w_inst = {w_inst:.2f} mm")
\end{lstlisting}

\begin{definition}
\textbf{numpy} (import\'e sous le nom \py{np}) est une biblioth\`eque externe
qui ajoute des capacit\'es de calcul scientifique \`a Python.
\py{np.arange(2.0, 6.5, 0.5)} g\'en\`ere un tableau de nombres de 2,0
\`a 6,0 (6,5 exclu) avec un pas de 0,5.
C'est l'\'equivalent Python d'une colonne de tableau dans Excel.
\end{definition}

Voici le code r\'eel du moteur de calcul~:

\begin{lstlisting}[style=python, caption={Extrait de \fichier{moteur.py} --- génération des portées}]
# L_min, pas, L_max viennent de la configuration et du produit
L_min = float(config.L_min_m)
pas   = float(config.pas_longueur_m)
L_max = p_rep.L_max_m   # portée max du produit (ex : 6.0 m)

# Génération du tableau de toutes les portées à tester
longueurs_m = np.arange(L_min, L_max + pas / 2, pas)
#              ^--- début    ^--- fin (+marge)  ^--- pas

# Note : le "+ pas/2" est une précaution technique.
# np.arange peut parfois manquer la borne finale à cause des
# imprécisions des nombres décimaux (6.0 devient 5.999...).
# En ajoutant la moitié du pas, on s'assure de bien inclure L_max.
\end{lstlisting}

% ─────────────────────────────────────────────────────────────────────────────
\section{Les conditions~: prendre des d\'ecisions}
% ─────────────────────────────────────────────────────────────────────────────

\begin{lstlisting}[style=python, caption={Conditions \py{if/elif/else} --- statut d'une portée}]
taux_determinant = 0.87   # résultat d'un calcul EC5

# Cas simple : admis ou refusé
if taux_determinant <= 1.0:
    statut = "admis"    # la poutre résiste
else:
    statut = "refusé"   # la poutre est insuffisante

# Cas avec plusieurs branches (elif = "sinon si")
if taux_determinant <= 0.80:
    commentaire = "large marge de sécurité"
elif taux_determinant <= 1.00:
    commentaire = "marge de sécurité acceptable"
else:
    commentaire = "section insuffisante"
\end{lstlisting}

Dans le code r\'eel (\fichier{moteur.py}), la d\'ecision est
l\'eg\`erement plus nuanc\'ee~:

\begin{lstlisting}[style=python, caption={Extrait de \fichier{moteur.py} --- logique de statut}]
statut_usage = els.get("statut_usage", "ok")

if statut_usage == "rejeté_usage":
    statut = "rejeté_usage"   # usage non supporté (ex : plancher dalle)
elif taux_determinant <= 1.0:
    statut = "admis"
else:
    statut = "refusé"
\end{lstlisting}

% ─────────────────────────────────────────────────────────────────────────────
\section{Les dictionnaires~: des fiches produit}
% ─────────────────────────────────────────────────────────────────────────────

\begin{maitre}
Dans ton catalogue bois, une fiche produit contient le lib\'ell\'e,
les dimensions, la classe de r\'esistance\ldots{}
Comment tu organiserais \c{c}a en Python~?
\end{maitre}

\begin{apprenti}
Je mettrais tout dans une sorte de tableau avec des \'etiquettes~?
\end{apprenti}

\begin{maitre}
Exactement~: c'est un \textbf{dictionnaire} (\py{dict}).
\end{maitre}

\begin{lstlisting}[style=python, caption={Dictionnaire --- fiche produit bois}]
produit = {
    "id_produit":         "EPX-063-175-6000",
    "b_mm":               63.0,
    "h_mm":               175.0,
    "L_max_m":            6.0,
    "classe_resistance":  "C24",
    "famille":            "bois_massif",
    "disponible":         True,
}

# Lecture d'une valeur par sa clé
hauteur   = produit["h_mm"]       # -> 175.0
est_dispo = produit["disponible"] # -> True

# Vérification de l'existence d'une clé
if "fournisseur" in produit:
    print(produit["fournisseur"])
else:
    print("Pas de fournisseur renseigné")
\end{lstlisting}

Dans le code r\'eel, les fonctions de v\'erification EC5 retournent
des dictionnaires avec toutes les valeurs interm\'ediaires du calcul~:

\begin{lstlisting}[style=python, caption={Dictionnaire retourné par \py{calculer\_flexion()}}]
def calculer_flexion(...):
    # ... calculs ...
    return {
        "f_m_d_MPa":       ...,  # résistance de calcul
        "sigma_m_MPa":     ...,  # contrainte calculée
        "taux_flexion_ELU": ..., # ratio de sollicitation (doit être <= 1.0)
    }
\end{lstlisting}

% ─────────────────────────────────────────────────────────────────────────────
\section{Les classes~: mod\'eliser les types de poutres}
% ─────────────────────────────────────────────────────────────────────────────

C'est ici que Python devient vraiment puissant. Jusqu'ici, on a
manipul\'e des valeurs et des fonctions isol\'ees. Les \textbf{classes}
permettent de regrouper des donn\'ees \emph{et} des comportements dans
une m\^eme entit\'e.

\begin{maitre}
Pense \`a une panne de toiture et \`a une solive de plancher.
Elles ont toutes les deux une largeur, une hauteur, une port\'ee\ldots{}
Mais comment sont calcul\'ees leurs charges~?
\end{maitre}

\begin{apprenti}
Diff\'eremment~! La panne subit les charges en biais \`a cause de la
pente du toit, tandis que la solive est horizontale.
\end{apprenti}

\begin{maitre}
Voil\`a exactement pourquoi ce code utilise des classes.
Regarde comment est d\'efinie la classe abstraite de base.
\end{maitre}

\begin{lstlisting}[style=python, caption={Extrait de \fichier{types\_poutre/base.py} --- classe abstraite \py{TypePoutre}}]
from abc import ABC, abstractmethod   # ABC = Abstract Base Class
import numpy as np

class TypePoutre(ABC):
    """Classe abstraite représentant un type de poutre bois.

    Toute sous-classe DOIT implémenter charges_lineaires().
    INTERDIT : if/match sur le type dans elu.py / els.py.
    """

    @abstractmethod
    def charges_lineaires(self, config, materiau, longueurs_m, combi):
        """Calcule les charges linéaires (kN/m) pour chaque portée.

        Cette méthode n'a PAS de code ici : c'est un "contrat".
        Chaque sous-classe (Panne, Solive, etc.) doit l'implémenter.
        """
        ...   # les "..." signifient : pas d'implémentation ici

    def longueur_deversement_m(self, longueur_m, entraxe_antideversement_mm):
        """Longueur de déversement (méthode COMMUNE à tous les types).

        Règle (EF-024) :
            pas d'anti-déversement  -> L_ef = L
            L <= 2 * entraxe        -> L_ef = L / 2
            L >  2 * entraxe        -> L_ef = entraxe
        """
        if entraxe_antideversement_mm <= 0:
            return longueur_m
        a_m = entraxe_antideversement_mm / 1000.0   # mm -> m
        if longueur_m <= 2 * a_m:
            return longueur_m / 2.0
        return a_m
\end{lstlisting}

\begin{definition}
\textbf{ABC} signifie \emph{Abstract Base Class} (classe de base abstraite).
Une classe abstraite est un \emph{contrat}~: elle dit
\guillemotleft\,toute sous-classe doit impl\'ementer ces m\'ethodes\,\guillemotright,
sans fournir l'impl\'ementation elle-m\^eme.
Le d\'ecorateur \py{@abstractmethod} marque les m\'ethodes obligatoires.
\end{definition}

Ensuite, chaque type de poutre h\'erite de cette classe et impl\'emente
sa propre logique de charges. Voici la classe \py{Panne} simplifi\'ee~:

\begin{lstlisting}[style=python, caption={Classe \py{Panne} --- implémentation pour toiture inclinée (simplifié)}]
import math
import numpy as np

class Panne(TypePoutre):
    """Panne de toiture inclinée.

    Hérite de TypePoutre et implémente charges_lineaires()
    en tenant compte de la pente.
    """

    def charges_lineaires(self, config, materiau, longueurs_m, combi):
        pente_rad = math.radians(config.pente_deg)
        entraxe_m = config.entraxe_m

        # Charges surfaciques (kN/m²) -> charges linéaires (kN/m)
        g_kNm2 = config.g_k_kNm2 + materiau.poids_propre_kNm2  # permanent
        s_kNm2 = config.s_k_kNm2                                # neige

        # Projection : les charges permanentes s'exercent sur la surface
        # inclinée, donc on les divise par cos(alpha) pour ramener
        # à la surface projetée horizontale
        g_inclinee = g_kNm2 / math.cos(pente_rad)

        q_G = g_inclinee * entraxe_m   # kN/m (par mètre de portée)
        q_S = s_kNm2 * entraxe_m       # kN/m (neige = charge verticale)

        # Combinaison EN 1990 (coefficients venant de "combi")
        q_d = combi.gamma_G * q_G + combi.gamma_Q1 * q_S

        # Moment et effort tranchant (poutre simplement appuyée)
        M_d = q_d * longueurs_m**2 / 8.0    # kN.m
        V_d = q_d * longueurs_m / 2.0       # kN

        return {
            "q_G_kNm": q_G * np.ones_like(longueurs_m),
            "q_S_kNm": q_S * np.ones_like(longueurs_m),
            "q_d_kNm": q_d * np.ones_like(longueurs_m),
            "M_d_kNm": M_d,
            "V_d_kN":  V_d,
        }

    def decomposer(self, charges_kNm, pente_rad):
        q_y = charges_kNm * math.cos(pente_rad)  # axe fort
        q_z = charges_kNm * math.sin(pente_rad)  # axe faible
        return q_y, q_z
\end{lstlisting}

\begin{important}
C'est le \textbf{polymorphisme}~: les modules de calcul EC5
(\fichier{elu.py}, \fichier{els.py}) appellent toujours
\py{type\_poutre.charges\_lineaires()} sans savoir si c'est une Panne,
une Solive ou un Chevron. Chaque sous-classe r\'epond \`a sa fa\c{c}on.
C'est \'el\'egant et extensible~: ajouter un nouveau type de poutre
ne n\'ecessite de modifier \emph{que} la nouvelle classe, sans toucher
aux calculs EC5.
\end{important}

% ─────────────────────────────────────────────────────────────────────────────
\section{Les modules et les imports}
% ─────────────────────────────────────────────────────────────────────────────

\begin{maitre}
On voit souvent \py{import numpy as np}, \py{import pandas as pd}\ldots{}
C'est quoi ces imports~?
\end{maitre}

\begin{apprenti}
Ce sont des biblioth\`eques ext\'erieures qu'on charge~?
\end{apprenti}

\begin{maitre}
Exactement. Imagine que tu travailles sur un chantier.
Tu n'emportes pas tous tes outils chaque matin~---
tu charges seulement ce dont tu as besoin pour la journ\'ee.
Les imports, c'est pareil.
\end{maitre}

\begin{lstlisting}[style=python, caption={Imports typiques dans ce projet}]
import numpy as np          # calcul numérique (tableaux, maths)
import pandas as pd         # manipulation de tableaux de données CSV
from pathlib import Path    # gestion propre des chemins de fichiers
import math                 # fonctions mathématiques (sin, cos, sqrt...)
from loguru import logger   # affichage de messages de log

# Import depuis un sous-module du projet lui-même
from abac_charpente.ec5.elu import verifier_elu
from abac_charpente.ec0.combinaisons import generer_combinaisons
\end{lstlisting}

Les biblioth\`eques cl\'es de ce projet sont~:

\begin{center}
\begin{tabular}{ll}
\toprule
\textbf{Biblioth\`eque} & \textbf{R\^ole}\\
\midrule
\py{numpy}    & Calcul num\'erique vectoris\'e (tableaux, maths)\\
\py{pandas}   & Lecture/\'ecriture CSV, manipulation de donn\'ees\\
\py{click}    & Interface ligne de commande\\
\py{pydantic} & Validation automatique des donn\'ees de configuration\\
\py{loguru}   & Affichage de messages (info, avertissement, erreur)\\
\py{matplotlib} & G\'en\'eration de graphiques (module \py{abac\_visuel})\\
\bottomrule
\end{tabular}
\end{center}

% ===========================================================================
\chapter{L'architecture du projet}
% ===========================================================================

% ─────────────────────────────────────────────────────────────────────────────
\section{Vue d'ensemble~: trois modules ind\'ependants}
% ─────────────────────────────────────────────────────────────────────────────

Le projet est d\'ecoup\'e en trois sous-packages ind\'ependants~:

\begin{center}
\begin{tikzpicture}[
    box/.style={
        draw, rounded corners=4pt,
        fill=bleuMaitreLight,
        minimum width=4.2cm, minimum height=1.4cm,
        align=center, font=\small\bfseries
    },
    arrow/.style={->, thick, color=bleuMaitre},
    label/.style={font=\scriptsize\itshape, color=gray}
]
\node[box] (sapeg) at (0,0)
    {\py{sapeg\_regen\_stock}\\[3pt]\small Ingestion du stock};

\node[box] (abac) at (5.5,0)
    {\py{abac\_charpente}\\[3pt]\small Calcul EC5};

\node[box] (visuel) at (11,0)
    {\py{abac\_visuel}\\[3pt]\small Graphiques};

\draw[arrow] (sapeg) -- (abac)
    node[midway, above, label] {stock\_enrichi.csv};
\draw[arrow] (abac) -- (visuel)
    node[midway, above, label] {portees\_admissibles.csv};
\end{tikzpicture}
\end{center}

\vspace{0.5em}

Cette s\'eparation suit un principe fondamental en informatique~:
chaque module a une \textbf{responsabilit\'e unique}.
Comme en charpente, on s\'epare les corps d'\'etat~---
le charpentier, le couvreur, le plaquiste~---
chacun fait son travail de fa\c{c}on autonome.

% ─────────────────────────────────────────────────────────────────────────────
\section{Structure des fichiers}
% ─────────────────────────────────────────────────────────────────────────────

\begin{lstlisting}[style=toml, language={}, caption={Arborescence du projet (simplifiée)}]
CAB_Chapente/
│
├── config.toml               <- Configuration principale
├── configs_calcul.toml       <- Scénarios de calcul EC5
├── configs_filtre.toml       <- Règles de filtrage du stock
├── ALL_PRODUIT_*.csv         <- Stock bois (export SAPEG)
│
├── src/
│   ├── abac_charpente/
│   │   ├── cli.py            <- Point d'entrée (ligne de commande)
│   │   ├── moteur.py         <- Orchestrateur principal
│   │   ├── config.py         <- Lecture de la configuration TOML
│   │   ├── registre.py       <- Cache de calculs (registre incrémental)
│   │   ├── ec0/
│   │   │   └── combinaisons.py  <- EN 1990 : combinaisons d'actions
│   │   ├── ec1/
│   │   │   ├── neige.py      <- EN 1991 : charges de neige
│   │   │   └── vent.py       <- EN 1991 : charges de vent
│   │   ├── ec5/
│   │   │   ├── elu.py        <- Etats Limites Ultimes (§6.1.x, §6.3.x)
│   │   │   ├── els.py        <- Etats Limites de Service (§7.2)
│   │   │   ├── double_flexion.py  <- Flexion biaxiale (§6.1.6)
│   │   │   ├── proprietes.py <- Tables kmod, kdef, gamma_M, EN 338
│   │   │   └── types_poutre/ <- Panne, Solive, Chevron, Sommier
│   │   └── data/             <- Tables normatives (CSV)
│   │       ├── materiaux_bois.csv    <- EN 338 : propriétés des classes
│   │       ├── kmod.csv              <- EC5 Tableau 3.1
│   │       ├── kdef.csv              <- EC5 Tableau 3.2
│   │       ├── gamma_m.csv           <- EC5 §2.4
│   │       └── limites_fleche_ec5.csv <- EC5 §7.2
│   │
│   ├── sapeg_regen_stock/    <- Pipeline stock SAPEG
│   └── abac_visuel/          <- Génération de graphiques
│
└── resultats/                <- Sorties (créées automatiquement)
    ├── portees_admissibles.csv
    ├── stock_enrichi.csv
    └── registre_calcul.csv
\end{lstlisting}

% ─────────────────────────────────────────────────────────────────────────────
\section{Le point d'entr\'ee~: la ligne de commande}
% ─────────────────────────────────────────────────────────────────────────────

Le programme se lance avec~:

\begin{lstlisting}[style=toml, language={}, caption={Lancement du logiciel}]
uv run abac calculer --config config.toml
\end{lstlisting}

\begin{maitre}
Comment Python sait-il quoi faire quand on tape cette commande~?
\end{maitre}

\begin{apprenti}
Il doit y avoir un fichier qui dit~:
\guillemotleft\,si on appelle la commande \py{abac}, d\'emarre ici\,\guillemotright~?
\end{apprenti}

\begin{maitre}
Oui. Ce fichier s'appelle \fichier{cli.py} et utilise la biblioth\`eque
\py{Click}. Les d\'ecorateurs \py{@click.command} et \py{@click.option}
transforment une fonction Python ordinaire en commande de terminal.
\end{maitre}

\begin{lstlisting}[style=python, caption={Extrait de \fichier{cli.py} --- point d'entrée simplifié}]
import click

@click.group()
def abac():
    """ABAC-Charpente : calcul de portées admissibles EC5."""
    pass

@abac.command("calculer")
@click.option("--config", default="config.toml",
              help="Fichier de configuration principal")
@click.option("--recalcul-complet", is_flag=True, default=False,
              help="Forcer le recalcul (ignorer le registre)")
def calculer(config, recalcul_complet):
    """Lance le calcul complet."""
    app_config = charger_config(config)
    moteur.lancer_calcul(app_config, recalcul_complet=recalcul_complet)
\end{lstlisting}

\begin{definition}
Un \textbf{d\'ecorateur} (le \py{@} devant une fonction) est une fa\c{c}on
d'enrober une fonction existante avec du comportement suppl\'ementaire,
sans modifier son code interne. Ici, \py{@click.command("calculer")} dit
\`a Click~: \guillemotleft\,la fonction suivante r\'epond \`a la
sous-commande \py{calculer}\,\guillemotright.
\end{definition}

% ─────────────────────────────────────────────────────────────────────────────
\section{Les mod\`eles de donn\'ees~: Pydantic}
% ─────────────────────────────────────────────────────────────────────────────

\begin{maitre}
Imagine que quelqu'un \'ecrit \py{g\_k\_kNm2 = "quarante"} dans le fichier
de configuration au lieu de \py{40.0}. Comment le programme sait-il que
c'est une erreur~?
\end{maitre}

\begin{apprenti}
Il faut qu'il v\'erifie les types des donn\'ees \`a la lecture~?
\end{apprenti}

\begin{maitre}
Exactement. C'est le r\^ole de \textbf{Pydantic}. Chaque configuration
est d\'ecrite par un mod\`ele de donn\'ees qui d\'efinit les types
attendus. Si un fichier TOML contient une valeur du mauvais type,
Pydantic l\`eve une erreur claire avant m\^eme que les calculs commencent.
\end{maitre}

\begin{lstlisting}[style=python, caption={\fichier{modeles/config\_calcul.py} --- modèle Pydantic (simplifié)}]
from pydantic import BaseModel

class ConfigCalcul(BaseModel):
    """Configuration d'un scénario de calcul EC5."""

    id_config_calcul: str          # identifiant unique (ex : "PANNE_INACC")
    type_poutre: str               # "Panne", "Solive", "Chevron", "Sommier"
    usage: str                     # "TOITURE_INACC", "PLANCHER_HAB", etc.

    L_min_m: float = 2.0           # portée minimale (m)
    pas_longueur_m: float = 0.5    # pas de calcul (m)

    pente_deg: float = 0.0         # pente toiture (degrés)
    entraxe_m: float = 1.0         # entraxe des poutres (m)
    classe_service: int = 1        # classe de service EC5 (1, 2 ou 3)

    g_k_kNm2: float                # charge permanente hors poids propre
    s_k_kNm2: float = 0.0          # neige (kN/m²)
    w_k_kNm2: float = 0.0          # vent (kN/m²)
    q_k_kNm2: float = 0.0          # exploitation (kN/m²)

    double_flexion: bool = False
    second_oeuvre: bool = False
    marge_securite: float = 0.0    # marge supplémentaire (0 à 1)
\end{lstlisting}

% ===========================================================================
\chapter{Les donn\'ees d'entr\'ee}
% ===========================================================================

% ─────────────────────────────────────────────────────────────────────────────
\section{Le stock bois~: le fichier ALL\_PRODUIT}
% ─────────────────────────────────────────────────────────────────────────────

Le stock est fourni par le logiciel SAPEG sous forme d'un fichier CSV
(valeurs s\'epar\'ees par des \py{|}). Chaque ligne est un produit disponible~:

\begin{lstlisting}[style=toml, language={}, caption={Exemple de lignes du fichier ALL\_PRODUIT}]
id_produit|libelle|b_mm|h_mm|L_max_m|classe_resistance
EPX-063-175|Epicéa C24 63x175|63.0|175.0|6.0|C24
EPX-075-225|Epicéa C24 75x225|75.0|225.0|7.0|C24
GL24H-160-320|GL24h 160x320|160.0|320.0|12.0|GL24h
\end{lstlisting}

\begin{maitre}
Pourquoi appelle-t-on ce format \guillemotleft\,CSV\,\guillemotright{}
si le s\'eparateur est \py{|} et non une virgule~?
\end{maitre}

\begin{apprenti}
C'est une bonne question\ldots{} CSV veut dire
\emph{Comma-Separated Values} mais en pratique on utilise
n'importe quel s\'eparateur~?
\end{apprenti}

\begin{maitre}
Oui, le terme CSV est devenu g\'en\'erique. SAPEG utilise le pipe
(\py{|}) car les lib\'ell\'es de produits peuvent contenir des virgules
ou des points-virgules. C'est un choix pratique pour \'eviter les
ambigu\"it\'es lors de la lecture.
\end{maitre}

Le code lit ce CSV avec Pandas~:

\begin{lstlisting}[style=python, caption={Lecture du stock CSV --- extrait de \fichier{sapeg\_regen\_stock/chargeur.py}}]
import pandas as pd

def charger_stock(chemin_csv: str) -> pd.DataFrame:
    """Charge le fichier ALL_PRODUIT CSV exporté par SAPEG."""
    df = pd.read_csv(
        chemin_csv,
        sep="|",            # séparateur pipe (pas une virgule !)
        encoding="latin-1"  # encodage SAPEG (attention aux accents !)
    )
    return df
\end{lstlisting}

\begin{definition}
Un \textbf{DataFrame Pandas} est l'\'equivalent Python d'une feuille Excel~:
un tableau avec des noms de colonnes, o\`u chaque ligne est un enregistrement.
On peut filtrer, trier, calculer sur des colonnes enti\`eres en une seule
instruction. Contrairement \`a une boucle Python, ces op\'erations sont
tr\`es rapides car Pandas les ex\'ecute en C interne.
\end{definition}

% ─────────────────────────────────────────────────────────────────────────────
\section{La configuration TOML}
% ─────────────────────────────────────────────────────────────────────────────

\begin{lstlisting}[style=toml, caption={config.toml --- configuration principale}]
[stock]
repertoire = "."              # répertoire contenant ALL_PRODUIT_*.csv
filtre_calcul = "charpente"   # nom du filtre stock à utiliser

[sortie]
fichier_csv    = "resultats/portees_admissibles.csv"
registre       = "resultats/registre_calcul.csv"
stock_enrichi  = "resultats/stock_enrichi.csv"

[calcul]
fichier_configs_calcul = "configs_calcul.toml"

[calcul.defaults]
longueur_appui_mm = 50.0    # longueur d'appui par défaut (mm)
k_c90             = 1.0     # facteur compression perp. au fil
taux_cible_appui  = 0.80    # taux cible pour calcul longueur appui min
\end{lstlisting}

Le fichier \fichier{configs\_calcul.toml} d\'efinit les sc\'enarios de calcul~:

\begin{lstlisting}[style=toml, caption={configs\_calcul.toml --- scénario panne de toiture}]
[[config_calcul]]
id_config_calcul = "PANNE_TOITURE_INACC"
type_poutre      = "Panne"
usage            = "TOITURE_INACC"
L_min_m          = 2.0
pas_longueur_m   = 0.5
# Listes -> expansion cartésienne automatique
pente_deg        = [4, 15, 30, 35, 40, 45]   # 6 pentes
entraxe_m        = [1.2, 1.7]                 # 2 entraxes
g_k_kNm2         = [0.4, 0.5, 0.6, 0.7, 0.8] # 5 charges perm.
# Total : 6 x 2 x 5 = 60 configurations générées !
s_k_kNm2         = 0.36
classe_service   = 1
double_flexion   = true
marge_securite   = 0
\end{lstlisting}

\begin{maitre}
Tu remarques que \py{pente\_deg} est une liste \py{[4, 15, 30, 35, 40, 45]}.
Qu'est-ce que l'\guillemotleft\,expansion cart\'esienne\,\guillemotright~?
\end{maitre}

\begin{apprenti}
Le programme cr\'ee une configuration distincte pour chaque combinaison
de pente, d'entraxe et de charge permanente~?
\end{apprenti}

\begin{maitre}
Exactement. Avec 6 pentes \texttimes{} 2 entraxes \texttimes{} 5 charges
permanentes, on obtient $6 \times 2 \times 5 = 60$ configurations \`a
partir d'une seule entr\'ee TOML. C'est un gain de temps consid\'erable.
\end{maitre}

% ─────────────────────────────────────────────────────────────────────────────
\section{Les tables normatives~: les CSV EC5}
% ─────────────────────────────────────────────────────────────────────────────

Le dossier \fichier{data/} contient les tables normatives de l'Eurocode~5
sous forme CSV, charg\'ees une seule fois au d\'emarrage~:

\begin{center}
\begin{tabular}{lll}
\toprule
\textbf{Fichier} & \textbf{Source} & \textbf{Contenu}\\
\midrule
\fichier{materiaux\_bois.csv} & EN~338 & Classes C16--C50, GL24--GL32\\
\fichier{kmod.csv}     & EC5 Tab.~3.1 & Facteurs $k_{mod}$ par famille/classe/dur\'ee\\
\fichier{kdef.csv}     & EC5 Tab.~3.2 & Facteurs de fluage $k_{def}$\\
\fichier{gamma\_m.csv} & EC5 \S{}2.4  & Coefficients partiels mat\'eriau $\gamma_M$\\
\fichier{limites\_fleche\_ec5.csv} & EC5 \S{}7.2 & Limites de fl\`eche par usage\\
\bottomrule
\end{tabular}
\end{center}

\begin{lstlisting}[style=toml, language={}, caption={Extrait de \fichier{materiaux\_bois.csv} (EN 338)}]
classe;famille;f_m_k_MPa;f_v_k_MPa;f_c90_k_MPa;E_0_mean_MPa;E_0_05_MPa
C16;bois_massif;16.0;1.8;2.2;8000;5400
C24;bois_massif;24.0;2.5;2.5;11000;7400
C30;bois_massif;30.0;3.0;2.7;12000;8000
GL24h;bois_lamelle_colle;24.0;2.7;2.5;11500;9600
\end{lstlisting}

\begin{maitre}
Pourquoi stocker ces tables dans des fichiers CSV plut\^ot que les coder
directement dans le programme~?
\end{maitre}

\begin{apprenti}
Parce que si la norme est mise \`a jour, on n'a qu'\`a modifier le fichier
CSV sans retoucher au code Python~?
\end{apprenti}

\begin{maitre}
Exactement. C'est la s\'eparation entre le \emph{code} (la logique) et
les \emph{donn\'ees} (les valeurs). Un principe fondamental de conception
logicielle.
\end{maitre}

% ===========================================================================
\chapter{Le moteur de calcul EC5}
% ===========================================================================

C'est le c\oe{}ur du logiciel. Ce chapitre explique comment le programme
passe d'un stock de poutres \`a des port\'ees admissibles v\'erifi\'ees
selon l'Eurocode.

% ─────────────────────────────────────────────────────────────────────────────
\section{Vue d'ensemble~: la boucle principale}
% ─────────────────────────────────────────────────────────────────────────────

Voici le flux de calcul complet, de haut en bas~:

\begin{center}
\begin{tikzpicture}[
    etape/.style={
        draw, rounded corners=3pt,
        fill=bleuMaitreLight,
        minimum width=8cm, minimum height=0.75cm,
        align=center, font=\small
    },
    arrow/.style={->, thick, color=bleuMaitre}
]
\node[etape] (e1) at (0,0)     {1. Charger le stock filtr\'e (CSV)};
\node[etape] (e2) at (0,-1.0)  {2. Grouper les produits par \py{id\_config\_materiau}};
\node[etape] (e3) at (0,-2.0)  {3. Pour chaque groupe $\times$ configuration EC5};
\node[etape] (e4) at (0,-3.0)  {4. G\'en\'erer tableau de port\'ees [L\_min \ldots L\_max]};
\node[etape] (e5) at (0,-4.0)  {5. G\'en\'erer combinaisons EN 1990 (ELU, ELS)};
\node[etape] (e6) at (0,-5.0)  {6. V\'erifier ELU (\S{}6.1.5, 6.1.6, 6.1.7, 6.3.3)};
\node[etape] (e7) at (0,-6.0)  {7. V\'erifier ELS --- fl\`eches (\S{}7.2)};
\node[etape] (e8) at (0,-7.0)  {8. Statut~: admis / refus\'e / rejet\'e\_usage};
\node[etape] (e9) at (0,-8.0)  {9. \'Ecrire CSV de sortie + mettre \`a jour le registre};

\draw[arrow] (e1) -- (e2); \draw[arrow] (e2) -- (e3);
\draw[arrow] (e3) -- (e4); \draw[arrow] (e4) -- (e5);
\draw[arrow] (e5) -- (e6); \draw[arrow] (e6) -- (e7);
\draw[arrow] (e7) -- (e8); \draw[arrow] (e8) -- (e9);
\end{tikzpicture}
\end{center}

% ─────────────────────────────────────────────────────────────────────────────
\section{Le groupement par mat\'eriau~: \'eviter les calculs en double}
% ─────────────────────────────────────────────────────────────────────────────

\begin{maitre}
Si ton stock contient 50~poutres C24 de 63\texttimes{}175, dois-tu
les calculer 50 fois~?
\end{maitre}

\begin{apprenti}
Non, elles ont toutes les m\^emes propri\'et\'es m\'ecaniques~!
Il suffit de calculer une fois et de copier le r\'esultat aux 50~produits.
\end{apprenti}

\begin{maitre}
C'est exactement ce que fait le code.
\end{maitre}

\begin{lstlisting}[style=python, caption={Groupement par \py{id\_config\_materiau} --- extrait de \fichier{moteur.py}}]
# groupes : clé = id matériau, valeur = liste de produits identiques
groupes: dict[str, list[ProduitValide]] = {}

for produit in produits_valides:
    # setdefault : si la clé n'existe pas, crée une liste vide d'abord
    groupes.setdefault(produit.id_config_materiau, []).append(produit)

# Résultat :
# groupes = {
#   "C24-063-175-6000": [prod_1, prod_2, prod_3],  # 3 refs 63x175
#   "C24-075-225-7000": [prod_4],                   # 1 ref  75x225
# }

for id_mat, produits_groupe in groupes.items():
    p_rep = produits_groupe[0]     # représentant du groupe pour le calcul

    # ... calculs EC5 (une seule fois pour tout le groupe) ...

    # Réplication : copier le résultat à chaque produit du groupe
    for produit in produits_groupe:
        # dataclasses.replace() crée une copie avec un champ modifié
        r_produit = replace(r_base, id_produit=produit.id_produit)
        tous.append(r_produit)
\end{lstlisting}

% ─────────────────────────────────────────────────────────────────────────────
\section{EN 1990~: les combinaisons d'actions}
% ─────────────────────────────────────────────────────────────────────────────

\begin{maitre}
Avant de v\'erifier une poutre, il faut savoir sous quelles charges
on la v\'erifie. Qu'est-ce qui charge une panne de toiture~?
\end{maitre}

\begin{apprenti}
Le poids propre du bois, le poids de la couverture, la neige\ldots{}
et peut-\^etre le vent~?
\end{apprenti}

\begin{maitre}
Exactement. Ces charges s'appellent des \textbf{actions} en langage Eurocode.
Et la norme EN~1990 dit qu'on ne peut pas simplement les additionner~---
il faut construire des \textbf{combinaisons} avec des coefficients
de s\'ecurit\'e, parce que la probabilit\'e que toutes les charges
soient simultan\'ement \`a leur valeur maximale est tr\`es faible.
\end{maitre}

La formule g\'en\'erale de combinaison ELU (\'Etat Limite Ultime ---
pour la r\'esistance) selon l'AN France est~:

\begin{equation}
E_{d} = \underbrace{\gamma_G \cdot G_k}_{\text{permanent}}
      + \underbrace{\gamma_{Q,1} \cdot Q_{k,1}}_{\text{action principale}}
      + \underbrace{\sum_{i>1} \psi_{0,i} \cdot \gamma_{Q,i} \cdot Q_{k,i}}_{\text{actions d'accompagnement}}
\label{eq:elu}
\end{equation}

\noindent o\`u~:
\begin{itemize}
  \item $G_k$ = charge permanente caract\'eristique (kN/m\textsuperscript{2})
  \item $Q_{k,1}$ = action variable principale (neige, vent ou exploitation)
  \item $Q_{k,i}$ = actions variables d'accompagnement
  \item $\gamma_G = 1{,}35$ (charges permanentes d\'efavorables)
  \item $\gamma_Q = 1{,}50$ (charges variables)
  \item $\psi_0$ = coefficient de combinaison (r\'eduit les actions secondaires)
\end{itemize}

Le code g\'en\`ere plusieurs cas~: une fois avec la neige dominante
(\py{ELU\_S}), une fois avec le vent dominant (\py{ELU\_W}),
une fois charges permanentes seules (\py{ELU\_G}), etc.

\begin{lstlisting}[style=python, caption={Extrait de \fichier{ec0/combinaisons.py} --- coefficients AN France et génération}]
# Coefficients réglementaires AN France
GAMMA_G = 1.35    # charges permanentes défavorables
GAMMA_Q = 1.50    # charges variables

# Coefficients ψ₀ : réduction des actions d'accompagnement
PSI_0_S = 0.5     # neige (altitude <= 1000 m)
PSI_0_W = 0.6     # vent
PSI_0_Q = 0.7     # exploitation catégorie A/B (habitation)

# Durées de charge (pour le facteur k_mod EC5 Tableau 3.1)
# G      -> permanent
# neige  -> court_terme
# vent   -> instantane

# Génération de la combinaison ELU avec la neige comme action principale
if s_k > 0:   # seulement s'il y a une charge de neige configurée
    combinaisons.append(CombinaisonEC0(
        id_combinaison  = "ELU_S",
        type_combinaison= "ELU_STR",
        charge_principale="S",
        gamma_G         = GAMMA_G,     # 1.35
        gamma_Q1        = GAMMA_Q,     # 1.50 * neige
        psi_0_Q2        = PSI_0_W if w_k > 0 else 0.0,  # vent accompagne
        duree_charge    = "court_terme",
    ))

# Et de même pour ELU_W (vent principal), ELU_Q (exploitation principale)...
\end{lstlisting}

Pour les \textbf{ELS} (\'Etats Limites de Service --- pour les fl\`eches),
les m\^emes actions sont combin\'ees \emph{sans} coefficients partiels~:

\begin{equation}
E_{d,ELS,car} = G_k + Q_{k,1} + \sum_{i>1} \psi_{0,i} \cdot Q_{k,i}
\label{eq:els_car}
\end{equation}

% ─────────────────────────────────────────────────────────────────────────────
\section{Le facteur \texorpdfstring{$k_{mod}$}{kmod}~:
         l'effet de la dur\'ee de charge}
% ─────────────────────────────────────────────────────────────────────────────

\begin{maitre}
Pourquoi le bois sous une charge de courte dur\'ee (neige)
r\'esiste mieux que sous une charge permanente~?
\end{maitre}

\begin{apprenti}
\`A cause du fluage~? Le bois se d\'eforme progressivement
sous charge longue, ce qui affaiblit sa r\'esistance~?
\end{apprenti}

\begin{maitre}
Exactement. L'Eurocode~5 traduit cela par le facteur $k_{mod}$ du
Tableau~3.1~: plus la charge est de longue dur\'ee, plus $k_{mod}$
est faible, et donc plus la r\'esistance de calcul $f_{m,d}$ est diminu\'ee.
\end{maitre}

La r\'esistance de calcul $f_{m,d}$ est obtenue \`a partir de la
r\'esistance caract\'eristique $f_{m,k}$ (issue des tables EN~338) par~:

\begin{equation}
f_{m,d} = f_{m,k} \cdot \frac{k_{mod}}{\gamma_M}
\label{eq:fmd}
\end{equation}

\noindent avec $\gamma_M = 1{,}3$ pour le bois massif (EC5 \S{}2.4).

\begin{center}
\begin{tabular}{ll}
\toprule
\textbf{Dur\'ee de charge} & \boldmath{$k_{mod}$} (bois massif, CS~1)\\
\midrule
Permanente & 0,60\\
Long terme & 0,70\\
Moyen terme & 0,80\\
Court terme (neige, vent) & 0,90\\
Instantan\'ee & 1,10\\
\bottomrule
\end{tabular}
\end{center}

Exemple num\'erique pour un bois C24 sous neige (court terme, CS~1)~:
\[
f_{m,d} = 24{,}0 \times \frac{0{,}90}{1{,}3} = 16{,}6 \text{ MPa}
\]

\begin{lstlisting}[style=python, caption={Usage de k\_mod dans \fichier{ec5/elu.py}}]
for combi in combinaisons:
    # Lecture de k_mod depuis la table kmod.csv
    # selon (famille du bois, classe de service, durée de charge)
    k_mod   = get_kmod(famille, classe_service, combi.duree_charge)
    gamma_M = get_gamma_m(famille)   # 1.3 pour bois massif

    # Résistance de calcul en flexion (MPa)
    f_m_d = materiau.f_m_k_MPa * k_mod / gamma_M
    # Ex : 24.0 × 0.90 / 1.3 = 16.6 MPa (neige, CS1)
\end{lstlisting}

% ─────────────────────────────────────────────────────────────────────────────
\section{ELU~--- V\'erification en flexion (\S{}6.1.6)}
% ─────────────────────────────────────────────────────────────────────────────

\subsection{La physique}

Sous une charge $q_d$ uniform\'ement r\'epartie, une poutre simplement
appuy\'ee de port\'ee $L$ est soumise \`a un moment maximum en trav\'ee~:

\begin{equation}
M_{d} = \frac{q_d \cdot L^2}{8}
\label{eq:moment}
\end{equation}

Ce moment cr\'ee une contrainte de flexion aux fibres extr\^emes~:

\begin{equation}
\sigma_{m,d} = \frac{M_d}{W_y} = \frac{M_d}{b \cdot h^2 / 6}
\label{eq:sigma}
\end{equation}

La v\'erification consiste \`a s'assurer que cette contrainte ne d\'epasse
pas la r\'esistance de calcul $f_{m,d}$~:

\begin{equation}
\frac{\sigma_{m,d}}{f_{m,d}} \leq 1{,}0
\label{eq:taux_flex}
\end{equation}

Ce rapport est appel\'e le \textbf{taux de sollicitation}.
Si taux $\leq 1$~: la poutre est admise. Sinon~: elle est refus\'ee.

\subsection{Le code}

\begin{lstlisting}[style=python, caption={\fichier{ec5/elu.py} --- fonction \py{calculer\_flexion()}}]
def calculer_flexion(
    materiau: ConfigMatériau,
    M_d_kNm: np.ndarray,   # tableau : un moment par longueur testée
    k_mod: float,
    gamma_M: float,
) -> dict[str, np.ndarray]:
    """Flexion §6.1.6 : σ_m_d = M_d / W_y ; taux = σ_m_d / f_m_d."""

    W_cm3     = materiau.W_cm3        # module de résistance (cm³)
    f_m_k_MPa = materiau.f_m_k_MPa   # résistance caractéristique (MPa)

    # Résistance de calcul (MPa) — Équation (4.3)
    f_m_d_MPa = f_m_k_MPa * k_mod / gamma_M

    # Contrainte de flexion (MPa)
    # Conversion des unités :
    #   M [kN.m] × 1e3 = M [N.m]
    #   M [N.m]  × 1e3 = M [N.mm]     (× 1e6 total)
    #   W [cm³]  × 1e3 = W [mm³]
    #   σ = M [N.mm] / W [mm³]  = N/mm² = MPa
    # -> σ [MPa] = M_kNm × 1e6 / (W_cm3 × 1e3) = M_kNm × 1e3 / W_cm3
    # ou encore : M_kNm × 1e3 / W_cm3 / 10.0
    #            car 1 kNm → MPa avec W en cm³ = ×1000/×10 = ×100...
    # Vérification : M=10 kNm, W=1000 cm³ : σ = 10×1000/1000/10 = 1 MPa ✓
    sigma_m_MPa = M_d_kNm * 1e3 / W_cm3 / 10.0

    return {
        "f_m_d_MPa":      np.full_like(sigma_m_MPa, f_m_d_MPa),
        "sigma_m_MPa":    sigma_m_MPa,
        "taux_flexion_ELU": sigma_m_MPa / f_m_d_MPa,
    }
\end{lstlisting}

\begin{maitre}
Pourquoi \py{M\_d\_kNm} est un \py{np.ndarray} plut\^ot qu'un simple \py{float}~?
\end{maitre}

\begin{apprenti}
Parce qu'on calcule pour toutes les longueurs en m\^eme temps~?
\py{M\_d\_kNm} contient par exemple 9~valeurs pour les 9~port\'ees
[2,0\,m, 2,5\,m, \ldots, 6,0\,m]~?
\end{apprenti}

\begin{maitre}
Exactement. NumPy applique l'op\'eration \py{* 1e3 / W\_cm3 / 10.0}
sur toutes les valeurs simultan\'ement, sans boucle explicite.
C'est ce qu'on appelle la \textbf{vectorisation}~--- bien plus rapide
qu'une boucle Python classique.
\end{maitre}

\begin{important}
La vectorisation NumPy est l'un des secrets de performance de ce logiciel.
L\`a o\`u une boucle Python traiterait les port\'ees une par une,
NumPy les traite toutes en parall\`ele au niveau du processeur.
Sur un grand stock, cela repr\'esente un facteur~100 de gain de vitesse.
\end{important}

% ─────────────────────────────────────────────────────────────────────────────
\section{ELU~--- V\'erification au cisaillement (\S{}6.1.7)}
% ─────────────────────────────────────────────────────────────────────────────

L'effort tranchant maximum dans une poutre simplement appuy\'ee est~:

\begin{equation}
V_{d} = \frac{q_d \cdot L}{2}
\end{equation}

La contrainte de cisaillement au niveau de l'axe neutre est~:

\begin{equation}
\tau_d = \frac{3}{2} \cdot \frac{V_d}{A_{eff}} = \frac{3}{2} \cdot \frac{V_d}{b_{eff} \cdot h}
\label{eq:tau}
\end{equation}

o\`u $b_{eff} = k_{cr} \cdot b$ est la largeur effective, r\'eduite
par le facteur $k_{cr}$ qui tient compte du risque de fissuration
en cisaillement (EC5 \S{}6.1.7(2)).
En pratique $k_{cr} = 0{,}67$ pour le bois massif non trait\'e
en surface.

La v\'erification~:

\begin{equation}
\frac{\tau_d}{f_{v,d}} \leq 1{,}0
\end{equation}

\begin{lstlisting}[style=python, caption={\fichier{ec5/elu.py} --- fonction \py{calculer\_cisaillement()}}]
def calculer_cisaillement(
    materiau: ConfigMatériau,
    V_d_kN: np.ndarray,
    k_mod: float,
    gamma_M: float,
    k_cr: float,        # facteur de fissuration (0.67 pour bois massif)
) -> dict[str, np.ndarray]:
    """Cisaillement §6.1.7."""
    b_eff_mm  = k_cr * materiau.b_mm          # largeur effective (mm)
    A_eff_mm2 = b_eff_mm * materiau.h_mm      # aire effective (mm²)

    f_v_d_MPa = materiau.f_v_k_MPa * k_mod / gamma_M  # résistance cisaillement

    # Contrainte (MPa) :
    # V [kN] × 1000 [N/kN] / A [mm²] × 1.5  (facteur 3/2 pour section rect.)
    tau_d_MPa = 1.5 * V_d_kN * 1000.0 / A_eff_mm2

    return {
        "f_v_d_MPa":          np.full_like(tau_d_MPa, f_v_d_MPa),
        "tau_MPa":            tau_d_MPa,
        "taux_cisaillement_ELU": tau_d_MPa / f_v_d_MPa,
    }
\end{lstlisting}

% ─────────────────────────────────────────────────────────────────────────────
\section{ELU~--- V\'erification \`a l'appui (\S{}6.1.5)}
% ─────────────────────────────────────────────────────────────────────────────

\`A l'appui, le bois est \'ecras\'e \textbf{perpendiculairement au fil},
ce qui est sa direction de moindre r\'esistance. La contrainte est~:

\begin{equation}
\sigma_{c,90,d} = \frac{V_d}{A_{appui}} = \frac{V_d}{b \cdot l_{appui}}
\end{equation}

La v\'erification EC5 \S{}6.1.5~:

\begin{equation}
\frac{\sigma_{c,90,d}}{k_{c,90} \cdot f_{c,90,d}} \leq 1{,}0
\end{equation}

Le code calcule aussi en retour la \textbf{longueur d'appui minimale}
$l_{min}$ pour satisfaire un taux cible de 0,80~:

\begin{equation}
l_{min} = \left\lceil \frac{V_d \cdot 1000}{b \cdot \text{taux\_cible} \cdot k_{c,90} \cdot f_{c,90,d}} \right\rceil \text{ mm}
\label{eq:lmin}
\end{equation}

\begin{lstlisting}[style=python, caption={\fichier{ec5/elu.py} --- appui et longueur minimale d'appui}]
def calculer_appui(materiau, V_d_kN, longueur_appui_mm, k_c90,
                   k_mod, gamma_M) -> dict:
    """Appui §6.1.5 : σ_c90_d = V_d / A_appui."""
    b_mm = materiau.b_mm
    f_c90_d_MPa = materiau.f_c90_k_MPa * k_mod / gamma_M

    A_appui_mm2 = b_mm * longueur_appui_mm          # aire d'appui (mm²)
    sigma_c90_MPa = V_d_kN * 1000.0 / A_appui_mm2  # MPa

    return {
        "sigma_c90_MPa":  sigma_c90_MPa,
        "taux_appui_ELU": sigma_c90_MPa / (k_c90 * f_c90_d_MPa),
    }

def calculer_longueur_appui_min(materiau, V_d_kN, k_c90, taux_cible,
                                k_mod, gamma_M) -> float:
    """Longueur d'appui minimale (mm) pour satisfaire le taux cible."""
    b_mm = materiau.b_mm
    f_c90_d_MPa = materiau.f_c90_k_MPa * k_mod / gamma_M
    denom = b_mm * taux_cible * k_c90 * f_c90_d_MPa
    if denom <= 0:
        return 0.0
    # math.ceil : arrondi au mm supérieur
    return math.ceil(V_d_kN * 1000.0 / denom)
\end{lstlisting}

% ─────────────────────────────────────────────────────────────────────────────
\section{ELU~--- D\'eversement lat\'eral (\S{}6.3.3)}
% ─────────────────────────────────────────────────────────────────────────────

\begin{maitre}
Qu'est-ce que le d\'eversement d'une poutre en bois~?
\end{maitre}

\begin{apprenti}
C'est quand une poutre \'elanc\'ee fl\'echit lat\'eralement au lieu de
rester dans son plan de chargement~? Une poutre haute et \'etroite,
peu retenue lat\'eralement, peut basculer sur le c\^ot\'e~?
\end{apprenti}

\begin{maitre}
Exactement. Plus la poutre est haute, \'etroite et longue (ou peu retenue),
plus elle risque de se d\'evier. L'EC5 \S{}6.3.3 quantifie ce risque
via le coefficient $k_{crit}$.
\end{maitre}

Le calcul part de la contrainte critique de d\'eversement
(EC5 \'Eq.~6.32) pour une section rectangulaire~:

\begin{equation}
\sigma_{m,crit} = \frac{0{,}78 \cdot b^2 \cdot E_{0,05}}{h \cdot L_{ef}}
\label{eq:sigmacrit}
\end{equation}

On en d\'eduit l'\'elancement relatif~:

\begin{equation}
\lambda_{rel,m} = \sqrt{\frac{f_{m,k}}{\sigma_{m,crit}}}
\label{eq:lambda}
\end{equation}

Et le coefficient $k_{crit}$ selon les seuils EC5 \S{}6.3.3(3)~:

\begin{equation}
k_{crit} = \begin{cases}
1{,}0 & \text{si } \lambda_{rel,m} \leq 0{,}75 \\
1{,}56 - 0{,}75 \, \lambda_{rel,m} & \text{si } 0{,}75 < \lambda_{rel,m} \leq 1{,}4 \\
1/\lambda_{rel,m}^2 & \text{si } \lambda_{rel,m} > 1{,}4
\end{cases}
\label{eq:kcrit}
\end{equation}

La v\'erification finale est~:

\begin{equation}
\frac{\sigma_{m,d}}{k_{crit} \cdot f_{m,d}} \leq 1{,}0
\label{eq:taux_dever}
\end{equation}

\begin{lstlisting}[style=python, caption={\fichier{ec5/elu.py} --- coefficient de déversement \py{calculer\_k\_crit()}}]
def calculer_k_crit(materiau, L_deversement_m) -> float:
    """Coefficient de déversement k_crit §6.3.3."""
    b_mm      = materiau.b_mm
    h_mm      = materiau.h_mm
    f_m_k_MPa = materiau.f_m_k_MPa
    E_005_MPa = materiau.E_0_05_MPa   # E_0,05 = module caractéristique

    L_dev_mm = L_deversement_m * 1000.0   # m -> mm

    # Contrainte critique de déversement (MPa) — EC5 Éq. (6.32)
    # σ_m,crit = 0.78 × b[mm]² × E_0,05[MPa] / (h[mm] × L_ef[mm])
    sigma_m_crit = 0.78 * b_mm**2 * E_005_MPa / (h_mm * L_dev_mm)

    # Élancement relatif
    lambda_rel_m = math.sqrt(f_m_k_MPa / sigma_m_crit)

    # Règles §6.3.3(3) : trois zones d'élancement
    if lambda_rel_m <= 0.75:
        return 1.0                            # pas de réduction
    elif lambda_rel_m <= 1.4:
        return 1.56 - 0.75 * lambda_rel_m    # réduction linéaire
    else:
        return 1.0 / lambda_rel_m**2          # réduction quadratique
\end{lstlisting}

\begin{maitre}
Quelle est la longueur de d\'eversement $L_{ef}$ utilis\'ee~?
\end{maitre}

\begin{apprenti}
La distance entre les \'el\'ements qui retiennent la poutre lat\'eralement~?
Comme les pannes secondaires, les liteaux ou les contrefiches~?
\end{apprenti}

\begin{maitre}
Exactement. C'est l'entraxe des anti-d\'eversements. Le code le g\`ere
dans la m\'ethode \py{longueur\_deversement\_m()} de la classe \py{TypePoutre}~:
\begin{itemize}
  \item Si pas d'anti-d\'eversement~: $L_{ef} = L$
  \item Si $L \leq 2 \times a$~: $L_{ef} = L/2$
  \item Si $L > 2 \times a$~: $L_{ef} = a$ (l'entraxe)
\end{itemize}
o\`u $a$ est l'entraxe des anti-d\'eversements.
\end{maitre}

% ─────────────────────────────────────────────────────────────────────────────
\section{ELS~--- Les fl\`eches (\S{}7.2)}
% ─────────────────────────────────────────────────────────────────────────────

\begin{maitre}
Apr\`es avoir v\'erifi\'e que la poutre ne casse pas (ELU),
que reste-t-il \`a v\'erifier~?
\end{maitre}

\begin{apprenti}
Qu'elle ne fl\'echit pas trop~? Une poutre qui tient mais qui fait
une grande fl\`eche visible, c'est probl\'ematique pour les rev\^etements
et d\'esagr\'eable visuellement.
\end{apprenti}

\begin{maitre}
Exactement. C'est l'objet des \'Etats Limites de Service (ELS).
Ici on ne v\'erifie plus la r\'esistance, mais la \emph{d\'eformation}.
\end{maitre}

\subsection{Fl\`eche instantan\'ee}

Pour une poutre simplement appuy\'ee sous charge uniforme~:

\begin{equation}
w_{inst} = \frac{5 \, q \, L^4}{384 \, E_{0,mean} \, I}
\label{eq:winst}
\end{equation}

\begin{lstlisting}[style=python, caption={\fichier{ec5/els.py} --- flèche instantanée}]
def calculer_w_inst(q_kNm, L_m, E_mean_MPa, I_cm4) -> float:
    """Flèche instantanée w_inst = 5qL⁴/(384EI) en mm."""
    # Conversions d'unités :
    q_Nmm = q_kNm * 1.0     # kN/m = N/mm (numériquement identique !)
    L_mm  = L_m * 1000.0    # m -> mm
    I_mm4 = I_cm4 * 1e4     # cm⁴ -> mm⁴

    return 5.0 * q_Nmm * L_mm**4 / (384.0 * E_mean_MPa * I_mm4)
\end{lstlisting}

\begin{maitre}
Pourquoi \py{q\_Nmm = q\_kNm * 1.0}~? La conversion est triviale~?
\end{maitre}

\begin{apprenti}
Oui~: 1 kN/m $=$ 1000 N / 1000 mm $=$ 1 N/mm, le facteur est 1.
\end{apprenti}

\begin{maitre}
Exact. La multiplication par 1{,}0 est gard\'ee pour la
\textbf{lisibilit\'e}~: le commentaire montre qu'il y a bien une conversion
d'unit\'e, m\^eme si le facteur num\'erique est 1. C'est le Principe~VI
du projet~: toujours rendre les unit\'es explicites.
\end{maitre}

\subsection{Fl\`eche finale~: le fluage}

Le bois flue sous charge permanente.
La fl\`eche finale tient compte de ce fluage via le facteur $k_{def}$~:

\begin{equation}
w_{fin} = w_{inst} \cdot (1 + k_{def})
\label{eq:wfin}
\end{equation}

\begin{lstlisting}[style=python, caption={\fichier{ec5/els.py} --- flèche finale avec fluage}]
def calculer_w_fin(w_inst, k_def) -> float:
    """Flèche finale w_fin = w_inst × (1 + k_def)."""
    return w_inst * (1.0 + k_def)

# Valeurs typiques de k_def (EC5 Tableau 3.2) :
# Bois massif, classe de service 1 : k_def = 0.6
# Bois massif, classe de service 2 : k_def = 0.8
# Bois lamellé-collé, CS1           : k_def = 0.6
\end{lstlisting}

\subsection{Limites de fl\`eche et taux ELS}

Les limites admissibles sont lues depuis \fichier{limites\_fleche\_ec5.csv}~:

\begin{center}
\begin{tabular}{lrrr}
\toprule
\textbf{Usage} & \boldmath{$w_{inst} \leq$} & \boldmath{$w_{fin} \leq$} & \boldmath{$w_2 \leq$}\\
\midrule
Toiture inaccessible & $L/150$ & $L/100$ & ---\\
Plancher habitation   & $L/300$ & $L/200$ & $L/300$\\
Plancher bureaux      & $L/300$ & $L/250$ & $L/350$\\
\bottomrule
\end{tabular}
\end{center}

Le taux ELS est calcul\'e de fa\c{c}on analogue au taux ELU~:

\begin{equation}
\text{taux}_{ELS,inst} = \frac{w_{inst}}{w_{inst,lim}}
\leq 1{,}0
\end{equation}

\begin{lstlisting}[style=python, caption={Extrait de \fichier{ec5/els.py} --- boucle de vérification ELS}]
for combi in combinaisons:
    if combi.type_combinaison not in ("ELS_CAR", "ELS_FREQ", "ELS_QPERM"):
        continue  # on ne calcule les flèches QUE pour les combi ELS

    # Charges via polymorphisme (comme pour l'ELU)
    charges = type_poutre.charges_lineaires(config, materiau, longueurs_m, combi)

    for i, L_m in enumerate(longueurs_m):
        q_d = float(charges["q_d_kNm"][i])

        # Limites depuis le CSV normatif
        limites = get_limites_fleche(config.usage, L_m)
        limite_inst_mm = limites["limite_inst_mm"]
        limite_fin_mm  = limites["limite_fin_mm"]

        # Calcul des flèches
        w_inst_mm  = calculer_w_inst(q_d, L_m, materiau.E_0_mean_MPa, materiau.I_cm4)
        w_creep_mm = k_def * w_inst_mm      # composante de fluage
        w_fin_mm   = calculer_w_fin(w_inst_mm, k_def)

        # Taux de sollicitation ELS
        taux_inst = w_inst_mm / limite_inst_mm
        taux_fin  = w_fin_mm  / limite_fin_mm
\end{lstlisting}

% ─────────────────────────────────────────────────────────────────────────────
\section{Double flexion (\S{}6.1.6)}
% ─────────────────────────────────────────────────────────────────────────────

\begin{maitre}
Pour une panne de toiture inclin\'ee, dans quel plan la charge
s'exerce-t-elle~?
\end{maitre}

\begin{apprenti}
La charge est verticale (poids, neige), mais la panne est inclin\'ee.
Donc la charge cr\'ee \`a la fois un moment dans le plan de la section
(axe fort $y$) \textbf{et} hors du plan (axe faible $z$)~?
\end{apprenti}

\begin{maitre}
Exactement. C'est la double flexion.
La d\'ecomposition de la charge $q_d$ selon les axes de la panne est~:
\end{maitre}

\begin{equation}
q_y = q_d \cdot \cos(\alpha) \qquad \text{(axe fort)}
\qquad
q_z = q_d \cdot \sin(\alpha) \qquad \text{(axe faible)}
\label{eq:decomp}
\end{equation}

\noindent o\`u $\alpha$ est la pente de la toiture. Les moments correspondants~:

\begin{equation}
M_{y,d} = \frac{q_y \cdot L^2}{8}, \qquad M_{z,d} = \frac{q_z \cdot L^2}{8}
\end{equation}

La v\'erification EN~1995 \S{}6.1.6 pour la double flexion donne
deux crit\`eres (les deux doivent \^etre satisfaits)~:

\begin{align}
\frac{\sigma_{m,y,d}}{f_{m,d}} + k_m \cdot \frac{\sigma_{m,z,d}}{f_{m,d}} &\leq 1{,}0
\label{eq:biaxial1}\\
k_m \cdot \frac{\sigma_{m,y,d}}{f_{m,d}} + \frac{\sigma_{m,z,d}}{f_{m,d}} &\leq 1{,}0
\label{eq:biaxial2}
\end{align}

\noindent avec $k_m = 0{,}7$ pour les sections rectangulaires.

\begin{lstlisting}[style=python, caption={Décomposition de charge et interaction biaxiale}]
import math

pente_rad = math.radians(config.pente_deg)

# Décomposition de la charge de calcul selon les axes de la panne
q_y = q_d * math.cos(pente_rad)   # composante axe fort
q_z = q_d * math.sin(pente_rad)   # composante axe faible

# Moments correspondants (poutre simplement appuyée)
M_y = q_y * L**2 / 8.0   # kN.m
M_z = q_z * L**2 / 8.0   # kN.m

# Contraintes (MPa) — même conversion que §6.1.6
sigma_y = M_y * 1e3 / materiau.W_cm3   / 10.0
sigma_z = M_z * 1e3 / materiau.W_z_cm3 / 10.0

# Deux taux d'interaction — §6.1.6 (les deux doivent être <= 1.0)
k_m = 0.7   # sections rectangulaires
taux_1 = sigma_y / f_m_d + k_m * sigma_z / f_m_d
taux_2 = k_m * sigma_y / f_m_d + sigma_z / f_m_d
\end{lstlisting}

\begin{maitre}
\`A quoi sert ce $k_m = 0{,}7$~?
\end{maitre}

\begin{apprenti}
Il r\'eduit la contribution de l'axe faible dans le crit\`ere ?
\end{apprenti}

\begin{maitre}
Oui, et inversement. C'est un facteur de r\'eduction de l'interaction~:
la norme reconnait que les deux moments n'atteignent pas leur valeur maximale
au m\^eme point de la section en m\^eme temps.
Pour une section carr\'ee ($b = h$), $k_m$ serait \'egal \`a 1{,}0
et les deux crit\`eres seraient identiques.
\end{maitre}

% ─────────────────────────────────────────────────────────────────────────────
\section{La marge de s\'ecurit\'e suppl\'ementaire}
% ─────────────────────────────────────────────────────────────────────────────

Le code permet d'ajouter une \textbf{marge de s\'ecurit\'e suppl\'ementaire}
via le param\`etre \py{marge\_securite} (entre 0 et 1).
Elle est appliqu\'ee \emph{apr\`es} tous les calculs EC5,
en multipliant tous les taux par $(1 + \text{marge})$~:

\begin{equation}
\text{taux\_effectif} = \text{taux}_{EC5} \times (1 + \text{marge\_securite})
\end{equation}

\begin{lstlisting}[style=python, caption={\fichier{moteur.py} --- application de la marge de sécurité (EF-026)}]
def _appliquer_marge(résultats: list[dict], marge: float) -> list[dict]:
    """Applique taux_effectif = taux × (1 + marge) sur tous les taux."""
    if marge == 0.0:
        return résultats   # optimisation : rien à faire si marge = 0

    champs_taux = [k for k in résultats[0].keys() if "taux" in k]
    for r in résultats:
        for k in champs_taux:
            v = r.get(k)
            if isinstance(v, float):
                r[k] = v * (1.0 + marge)
    return résultats
\end{lstlisting}

Exemple~: avec \py{marge\_securite = 0.15} et un taux calcul\'e de 0{,}85~:
\[
\text{taux\_effectif} = 0{,}85 \times 1{,}15 = 0{,}978 \leq 1{,}0
\quad \Rightarrow \text{ admis, mais conservateur}
\]

% ===========================================================================
\chapter{Les sorties}
% ===========================================================================

% ─────────────────────────────────────────────────────────────────────────────
\section{Le CSV de r\'esultats}
% ─────────────────────────────────────────────────────────────────────────────

Le fichier \fichier{resultats/portees\_admissibles.csv} est la sortie principale.
Chaque ligne correspond \`a une combinaison
(produit $\times$ config\_calcul $\times$ port\'ee $\times$ combinaison EN~1990).

Colonnes principales~:

\begin{center}
\begin{tabular}{ll}
\toprule
\textbf{Colonne} & \textbf{Description}\\
\midrule
\py{id\_produit} & R\'ef\'erence du produit bois\\
\py{id\_config\_calcul} & Identifiant du sc\'enario EC5\\
\py{type\_poutre} & Panne, Solive, Chevron, Sommier\\
\py{longueur\_m} & Port\'ee v\'erifi\'ee (m)\\
\py{statut} & \py{admis} / \py{refus\'e} / \py{rejet\'e\_usage}\\
\py{taux\_determinant} & Taux de sollicitation maximal\\
\py{verification\_determinante} & Quelle v\'erification est critique\\
\py{clause\_EC5} & R\'ef\'erence normative (\S{}6.1.6, \S{}7.2, etc.)\\
\py{sigma\_m\_MPa} & Contrainte de flexion calcul\'ee\\
\py{taux\_flexion\_ELU} & Taux ELU flexion\\
\py{tau\_MPa} & Contrainte de cisaillement calcul\'ee\\
\py{taux\_cisaillement\_ELU} & Taux ELU cisaillement\\
\py{taux\_appui\_ELU} & Taux ELU appui\\
\py{longueur\_appui\_min\_mm} & Longueur d'appui minimale calcul\'ee\\
\py{k\_crit} & Coefficient de d\'eversement\\
\py{taux\_deversement\_ELU} & Taux ELU d\'eversement\\
\py{w\_inst\_mm} & Fl\`eche instantan\'ee (mm)\\
\py{w\_fin\_mm} & Fl\`eche finale avec fluage (mm)\\
\py{taux\_ELS\_inst} & Taux ELS fl\`eche instantan\'ee\\
\py{taux\_ELS\_fin} & Taux ELS fl\`eche finale\\
\bottomrule
\end{tabular}
\end{center}

\begin{maitre}
Comment le programme sait-il quelle v\'erification est la plus critique~?
\end{maitre}

\begin{apprenti}
Il prend le taux le plus \'elev\'e parmi toutes les v\'erifications~?
\end{apprenti}

\begin{maitre}
Exactement~:
\end{maitre}

\begin{lstlisting}[style=python, caption={Extrait de \fichier{moteur.py} --- vérification déterminante}]
def _trouver_verification_determinante(elu, els, df, taux_max):
    """Trouve quelle vérification a le taux le plus élevé."""
    candidates = {
        "flexion_ELU":      elu.get("taux_flexion_ELU", 0),
        "cisaillement_ELU": elu.get("taux_cisaillement_ELU", 0),
        "appui_ELU":        elu.get("taux_appui_ELU", 0),
        "deversement_ELU":  elu.get("taux_deversement_ELU", 0),
        "fleche_inst":      els.get("taux_ELS_inst", 0),
        "fleche_fin":       els.get("taux_ELS_fin", 0),
    }
    # max() sur un dict : trouve la clé avec la valeur maximale
    return max(candidates, key=lambda k: candidates[k] or 0)

# Correspondance vérification -> clause normative
mapping_clause = {
    "flexion_ELU":      "§6.1.6",
    "cisaillement_ELU": "§6.1.7",
    "appui_ELU":        "§6.1.5",
    "deversement_ELU":  "§6.3.3",
    "fleche_inst":      "§7.2",
    "fleche_fin":       "§7.2",
}
\end{lstlisting}

% ─────────────────────────────────────────────────────────────────────────────
\section{Le registre incr\'emental}
% ─────────────────────────────────────────────────────────────────────────────

\begin{maitre}
Si le stock change l\'eg\`erement (ajout de deux produits),
faut-il recalculer toutes les port\'ees~?
\end{maitre}

\begin{apprenti}
Non, seulement pour les nouveaux produits.
Les produits existants n'ont pas chang\'e.
\end{apprenti}

\begin{maitre}
C'est exactement ce que fait le \textbf{registre de calcul}.
Il m\'emorise les paires (mat\'eriau, config) d\'ej\`a calcul\'ees
dans un fichier CSV persistant.
\end{maitre}

\begin{lstlisting}[style=python, caption={Registre incrémental --- extrait de \fichier{registre.py}}]
class RegistreCalcul:
    """Mémorise les calculs déjà effectués pour éviter les doublons."""

    def est_calcule(self, id_mat: str, id_calc: str) -> bool:
        """Retourne True si ce couple a déjà été calculé."""
        if self._force_recalcul:
            return False   # ignorer le registre si recalcul forcé
        return (id_mat, id_calc) in self._calcules

    def enregistrer(self, id_mat: str, id_calc: str, statut: str):
        """Marque ce couple comme calculé."""
        self._calcules.add((id_mat, id_calc))
\end{lstlisting}

Et dans la boucle principale du moteur~:

\begin{lstlisting}[style=python, caption={Usage du registre dans \fichier{moteur.py}}]
for config in app_config.configs_calcul:
    id_calc = config.id_config_calcul

    # Si déjà calculé -> sauter (sauf si --recalcul-complet)
    if not registre._force_recalcul and registre.est_calcule(id_mat, id_calc):
        if verbose:
            logger.info(f"Déjà calculé ({id_mat}, {id_calc}) - ignoré.")
        continue   # "continue" = passer à l'itération suivante

    # ... calculs EC5 ...

    registre.enregistrer(id_mat, id_calc, "calcule")
\end{lstlisting}

Pour forcer un recalcul complet~:

\begin{lstlisting}[style=toml, language={}, caption={Forcer le recalcul complet}]
uv run abac calculer --config config.toml --recalcul-complet
\end{lstlisting}

% ─────────────────────────────────────────────────────────────────────────────
\section{Les graphiques}
% ─────────────────────────────────────────────────────────────────────────────

Le module \py{abac\_visuel} lit le CSV de r\'esultats et g\'en\`ere
des graphiques de port\'ees admissibles avec Matplotlib~:

\begin{lstlisting}[style=toml, language={}, caption={Lancement de la génération graphique}]
uv run abac-visuel generer \
    --donnees resultats/portees_admissibles.csv \
    --configs configs_calcul.toml \
    --sortie  resultats/graphiques
\end{lstlisting}

Le module \fichier{abac\_visuel/generateur.py} lis le CSV,
groupe les r\'esultats par configuration de calcul,
et pour chaque configuration, trace les port\'ees admissibles
en fonction des dimensions de section (b, h), avec les taux de
sollicitation repr\'esent\'es par une \'echelle de couleurs.

% ===========================================================================
\chapter*{Conclusion~: le voyage complet}
\addcontentsline{toc}{chapter}{Conclusion}
% ===========================================================================

\begin{maitre}
R\'esumons notre voyage. Qu'est-ce que fait ce logiciel,
dans tes propres mots~?
\end{maitre}

\begin{apprenti}
Il lit un stock de poutres bois, g\'en\`ere toutes les configurations
de calcul (en faisant le produit cart\'esien des param\`etres),
et pour chaque poutre et chaque port\'ee, il v\'erifie la r\'esistance
(ELU~: flexion, cisaillement, appui, d\'eversement) et les d\'eformations
(ELS~: fl\`eches instantan\'ee et finale) selon l'Eurocode~5.
Il stocke les r\'esultats dans un CSV et m\'emorise ce qui a d\'ej\`a
\'et\'e calcul\'e pour ne pas tout recalculer \`a chaque fois.
\end{apprenti}

\begin{maitre}
Et Python dans tout \c{c}a~?
\end{maitre}

\begin{apprenti}
Python orchestre tout. Les \textbf{variables} stockent les donn\'ees
(avec l'unit\'e dans le nom). Les \textbf{fonctions} encapsulent les formules
(une fonction = une formule de la norme). Les \textbf{classes} mod\'elisent
les types de poutres avec leur comportement propre~: c'est le polymorphisme.
Les \textbf{boucles} parcourent le stock et les port\'ees. NumPy
\textbf{vectorise} les calculs pour la performance. Pandas lit et \'ecrit
les CSV. Et Pydantic valide les donn\'ees d'entr\'ee.
\end{apprenti}

\begin{important}
Les concepts Python rencontr\'es dans ce logiciel~--- variables, fonctions,
classes, polymorphisme, boucles, dictionnaires, modules~--- sont exactement
ceux qu'on enseigne en premi\`ere ann\'ee de programmation.
La diff\'erence, c'est qu'ici ils sont mis au service d'un probl\`eme
d'ing\'enierie r\'eel et norm\'e.
La structure du code \textbf{refl\`ete} la structure de la norme~:
un module \py{ec0/} pour EN~1990, un module \py{ec1/} pour EN~1991,
et un module \py{ec5/} pour EN~1995.
Ce n'est pas un hasard~: c'est une conception d\'elib\'er\'ee.
\end{important}

\vspace{2em}
\begin{center}
\rule{0.5\textwidth}{0.5pt}
\end{center}
\vspace{1em}

\begin{tcolorbox}[
  colback=bleuMaitreLight, colframe=bleuMaitre,
  width=0.9\textwidth, center
]
\centering\itshape
\guillemotleft\, La connaissance s'acqui\`ere par l'exp\'erience,\\
tout le reste n'est qu'information.\,\guillemotright\\[4pt]
{\normalsize --- Albert Einstein}
\end{tcolorbox}

\end{document}
